\subsection{Módulos fielmente planos}
\label{subsection:0.6.4-fr}

\begin{env}[6.4.1]
\label{0.6.4.1-fr}
Para un $A$-módulo $M$, las cuatro propiedades siguientes son equivalentes:
\begin{enumerate}
  \item[$a)$] Para que una sucesión $N'\to N\to N''$ de $A$-módulos sea exacta,
    es necesario y suficiente que la sucesión
    $M\otimes N'\to M\otimes N\to M\otimes N''$ sea exacta;
  \item[$b)$] $M$ es plano y para todo $A$-módulo $N$, la relación
    $M\otimes N=0$ implica $N=0$;
  \item[$c)$] $M$ es plano y para todo homomorfismo $v:N\to N'$ de
    $A$-módulos, la relación $1_M\otimes v=0$ implica $v=0$, siendo $1_M$
    el automorfismo identidad de $M$;
  \item[$d)$] $M$ es plano y para todo ideal maximal $\mathfrak{m}$ de
    $A$, $\mathfrak{m}M\neq M$.
\end{enumerate}

Cuando $M$ verifica estas condiciones, se dice que $M$ es un $A$-módulo
\emph{fielmente plano}; $M$ es entonces necesariamente un módulo
\emph{fiel}. Además, si $u:N\to N'$ es un homomorfismo de $A$-módulos,
para que $u$ sea inyectivo (resp. sobreyectivo, biyectivo), es necesario y suficiente
que $1\otimes u:M\otimes N\to M\otimes N'$ lo sea.
\end{env}

\oldpage[0\textsubscript{I}]{58}

\begin{env}[6.4.2]
\label{0.6.4.2-fr}
Un módulo libre $\neq\{0\}$ es fielmente plano; lo mismo ocurre con la
suma directa de un módulo plano y de un módulo fielmente plano. Si $S$ es
una parte multiplicativa de $A$, $S^{-1}A$ es un $A$-módulo fielmente
plano solo si $S$ está formado por elementos invertibles (por tanto $S^{-1}A=A$).
\end{env}

\begin{env}[6.4.3]
\label{0.6.4.3-fr}
Sea
\[
  0\longrightarrow M'\longrightarrow M\longrightarrow M''\longrightarrow 0
\]
una sucesión exacta de $A$-módulos; si $M'$ y $M''$ son planos, y si uno
de ellos es fielmente plano, entonces $M$ es fielmente plano.
\end{env}

\begin{env}[6.4.4]
\label{0.6.4.4-fr}
Sean $A$ y $B$ dos anillos, $M$ un $A$-módulo, $N$ un
$(A,B)$-bimódulo. Si $M$ es fielmente plano y si $N$ es un $B$-módulo
fielmente plano, entonces $M\otimes_A N$ es un $B$-módulo fielmente plano. En
particular, si $M$ y $N$ son dos $A$-módulos fielmente planos, lo mismo
ocurre con $M\otimes_A N$. Si $B$ es una $A$-álgebra y si $M$ es un
$A$-módulo fielmente plano, el $B$-módulo $M_{(B)}$ es fielmente plano.
\end{env}

\begin{env}[6.4.5]
\label{0.6.4.5-fr}
Si $M$ es un $A$-módulo fielmente plano y si $S$ es una parte
multiplicativa de $A$, $S^{-1}M$ es un $S^{-1}A$-módulo fielmente plano,
puesto que $S^{-1}M=M\otimes_A(S^{-1}A)$ (6.4.4). Recíprocamente, si para todo
ideal maximal $\mathfrak{m}$ de $A$, $M_\mathfrak{m}$ es un
$A_\mathfrak{m}$-módulo fielmente plano, $M$ es un $A$-módulo fielmente
plano, pues $M$ es $A$-plano (6.3.3), y se tiene
\[
  M_\mathfrak{m}/\mathfrak{m}M_\mathfrak{m}
  =(M\otimes_A A_\mathfrak{m})
  \otimes_{A_\mathfrak{m}}
  (A_\mathfrak{m}/\mathfrak{m}A_\mathfrak{m})
  =M\otimes_A(A/\mathfrak{m})
  =M/\mathfrak{m}M
\]
por tanto, la hipótesis implica que $M/\mathfrak{m}M\neq0$ para todo ideal
maximal $\mathfrak{m}$ de $A$, lo que prueba nuestra afirmación (6.4.1).
\end{env}

